% Block 8 of the slides (08_failures.tex) and the conclusion (90_conclusion.tex).
\section{Failure modes and responsibility}\label{sec:failures}

\onslides{the planted bug; ``plausible is not correct''; the bug your tests cannot
see; two rogues' galleries; you are the author; recording which tool was used.}

\subsection{The planted bug}

Participants apply a patch to their own repository with \cmd{git apply}, make the Rabi
plot, and run the tests. The patch changes one character: $\Omega/2 \to \Omega$ in
\cmd{hamiltonian}. Its effect at $\Delta = 0$: the Hamiltonian becomes $\Omega\sigma_x$, so
\begin{equation}
  \Pe(t) = \sin^2(\Omega t) \quad\text{instead of}\quad \sin^2(\Omega t/2).
\end{equation}
The plot still shows smooth, bounded, periodic oscillations between $0$ and $1$, at twice
the frequency. Without a reference, the eye accepts both curves (slide ``Plausible is not
correct''). The $\pi$-pulse test fails at once: at $t = \pi/\Omega$ the buggy population
is $\sin^2\pi = 0$ instead of $1$. Off resonance the bug changes amplitude and frequency
together, $\Pe = \Omega^2/(\Omega^2 + \Delta^2/4)\,\sin^2(\sqrt{\Omega^2 + \Delta^2/4}\,t)$,
so the analytic test fails at every $\Delta$. Find the bug with \cmd{git diff}; undo it with
\cmd{git restore .}.

\begin{placeholder}[Placeholder: \texttt{exercises/planted\_bug.patch}]
Quote the patch once D3 provides it (it must touch only the line of \texttt{hamiltonian},
so it applies after stages 2--4).
\end{placeholder}

\subsection{The bug your tests cannot see}

Flip the sign of the Hamiltonian term in Eq.~\eqref{eq:lindblad}: $-i[H, \rho] \to
+i[H, \rho]$. This is evolution under $-H$, equivalently backwards in time for the
coherent part. All populations in the course are unchanged, because
\begin{itemize}
  \item the Rabi formula~\eqref{eq:rabi} is even in $\Omega$ and in $\Delta$;
  \item the Ramsey and echo results~\eqref{eq:ramsey} depend on $\cos(\Delta\tau)$ only;
  \item the dissipators do not involve $H$.
\end{itemize}
So every test of stages 1--4 passes. Only phases change sign: at $\Delta = 0$,
$\langle\sigma_y\rangle(t) = \sin(\Omega t)$ (Section~\ref{sec:stage1}) becomes
$-\sin(\Omega t)$. The lesson on the slide: your tests decide which bugs you can see, so
test every quantity you will publish. A single test of $\langle\sigma_y\rangle$, or of
the sign of a Ramsey phase $\varphi$ with a known detuning, closes this gap.

\subsection{Rogues' gallery}

\begin{center}
\renewcommand{\arraystretch}{1.25}
\begin{tabularx}{\linewidth}{@{}p{0.36\linewidth}X@{}}
\toprule
\textbf{Failure} & \textbf{Why it happens, and the defence} \\
\midrule
Invented or misquoted reference & Plausible text with no lookup (Section~\ref{sec:llm}).
  Open the link; check that the paper exists \emph{and} says that. \\
Factor of 2, wrong sign & Conventions differ between sources ($\Omega$ vs.\ $\Omega/2$,
  $\Delta = \omega_q - \omega_d$ or the reverse). Analytic and limiting-case tests. \\
Test weakened to pass & Rewarded for passing checks (Section~\ref{sec:llm}). Commit tests
  first; read every test diff. \\
Convention silently changed & Fourier normalisation, $\hbar = 1$, basis order, sign of
  $\Delta$. Conventions in \cmd{AGENTS.md}, enforced by tests. \\
Invented or outdated function names & Knowledge cutoff, little Julia in the training data.
  Run it; pin versions in \cmd{Manifest.toml}. \\
Hz where rad/\textmu s was meant & A factor of $2\pi$ that looks harmless. Units in
  names and docstrings; a test with a known period. \\
Confident summary of a paper not read & The model writes what such a paper would
  plausibly say. Ask for page and equation numbers; check them. \\
``All tests pass'' that were never run & A summary is a claim like any other. Run them
  yourself; CI runs them anyway. Review the diff, not the summary. \\
\bottomrule
\end{tabularx}
\end{center}

\subsection{You are the author}

\textbf{Accountability.} An AI tool is not an author, and every claim in the paper is
yours. Nature's 2023 editorial \rf{nature-rules} set two ground rules that most
publishers now follow: no AI tool is credited as an author, and its use is documented in
the methods or acknowledgements. The American Physical Society updated its policy in
June 2026 \rf{aps-news,aps-policy}. It allows substantive uses (literature synthesis, data
analysis, scientific reasoning, figures) with disclosure, and states that an AI program
cannot be held accountable and therefore cannot be an author. The disclosure names the
tool and version, how it assisted, how it was directed, and how its output was verified;
AI used in the research goes into the methods section, AI-made figures into their
captions, other substantive uses into the acknowledgements. Light editing needs no
disclosure.

\textbf{Data.} When an agent reads a file, its content is sent to the provider. Unpublished
data, collaborators' code and personal data leave your machine. Check what your
institution and your collaborations allow, and what the provider's terms say about
retention and training. Referees take note: APS forbids uploading manuscripts under review
to tools whose terms allow training on them \rf{aps-policy}.

\textbf{Disclosure.} Read your journal's policy before you submit. The APS
requirement ``how its output was verified'' is easy to meet if the verification is in
the repository: the tests, the independent routes, the claim statuses.

\textbf{Reproducibility.} Record which tool and model helped, and when. A commit trailer
costs nothing and is searchable:
\begin{lstlisting}
Add Ramsey sequence with fringe-frequency test

Assisted-by: <agent> (<model>)
\end{lstlisting}
Some agents add a trailer themselves (for example \cmd{Co-Authored-By}); keep it. Git
treats any \cmd{Key: value} lines at the end of the message as trailers;
\cmd{git log --grep "Assisted-by"} finds them later.

\subsection{Conclusion of the course}

The six take-aways on the closing slide, with the section that explains each:
\begin{enumerate}
  \item An LLM is a very large fit: plausible is not the same as correct
        (Section~\ref{sec:llm}).
  \item An agent closes the loop, so it can check its work, if you give it checks
        (Section~\ref{sec:agents}).
  \item Context is finite: one task per session, write the rest down
        (Section~\ref{sec:context}).
  \item Git is the safety net: commit often, read every diff (Section~\ref{sec:git}).
  \item Structure, pinned environments and CI: every result traceable and checked
        (Section~\ref{sec:structure}).
  \item Verification is the job: tests that encode physics, written first, read by you
        (Sections~\ref{sec:verification} and~\ref{sec:physics}).
\end{enumerate}

\begin{qa}
\Q{Do I have to disclose that I used an agent to write analysis code?}
For APS journals, yes: AI used in conducting the research, such as data analysis, is
disclosed in the methods section \rf{aps-policy}. Other publishers differ; check.

\Q{Can I paste my unpublished data into the agent?}
Only if your institution, your collaboration and the provider's terms allow it. When in
doubt, use synthetic data with the same structure (stage~4 is an example).

\Q{The agent's summary says all tests pass. Is that enough?}
No. Run them yourself, or let CI run them. Review the diff, not the summary.
\end{qa}
